\section{What Surprised Us}

The headline result is that the binary heap was the fastest implementation on
every graph in both Dijkstra variants, and that the pairing heap never
overtook it. Hypothesis~1.2 predicted a crossover on dense graphs, where the
pairing heap's cheap decrease-key should dominate, yet on \code{uniform_dense},
with 50 edges per vertex, the binary heap still takes $1.69$~ms against
$2.43$~ms. We expected the explanation to be cache locality. The counters show
that locality helps, since the binary heap has the fewest L1 misses per
vertex, but the larger reason was one we had not anticipated at all.

\subsection{Decrease-key is rarely called}

Hypothesis~1.2 assumed that on dense graphs the $m - n$ possible decrease-key
calls would overwhelm the $n$ extractions. Counting the calls tells a
different story. The pointer heaps allocate exactly one node per push, so the
allocation counts in Section~4.4 give the number of pushes in each variant,
and the difference between the lazy and \code{decrease_key} counts is the
number of decrease-key calls (Table~\ref{tab:decrease-key-counts}).

\begin{table}[H]
  \centering
  \small
  \caption{Decrease-key calls compared with the $m - n$ assumed in
  Hypothesis~1.2.}
  \label{tab:decrease-key-counts}
  \begin{tabular}{@{}lrrr@{}}
    \toprule
    \textbf{Graph} & $\boldsymbol{m - n}$ & \textbf{Decrease-key calls} & \textbf{Per vertex} \\
    \midrule
    \code{smoke_1k}         & $3{,}000$   & $304$        & $0.30$ \\
    \code{uniform_sparse}   & $30{,}000$  & $3{,}066$    & $0.31$ \\
    \code{uniform_medium}   & $180{,}000$ & $18{,}764$   & $0.94$ \\
    \code{uniform_dense}    & $490{,}000$ & $22{,}814$   & $2.28$ \\
    \code{grid_224}         & $149{,}632$ & $16{,}029$   & $0.32$ \\
    \code{hub_skewed}       & $100{,}000$ & $3{,}362$    & $0.17$ \\
    \code{layered_frontier} & $138{,}400$ & $19{,}811$   & $0.99$ \\
    \bottomrule
  \end{tabular}
\end{table}

On \code{uniform_dense} only $22{,}814$ of the $490{,}000$ possible calls
happen, under $5\%$, and no graph reaches even $15\%$. Most relaxations fail
because, with random weights, the first tentative distance found for a vertex
is usually already close to its final distance. This matches a classical
result: for random edge weights the expected number of decrease-key
operations grows only like $n \log(m/n)$ (Noshita, 1985), about $39{,}000$ for
\code{uniform_dense}, the same order of magnitude as we measured.

This one fact explains three of our wrong predictions. The pairing heap's
cheaper decrease-key saves time on only about two operations per vertex, too
few to pay for its pointer-based layout, so there is no crossover. The lazy
variant never holds anywhere near $m$ entries, since it makes at most $3.3$
pushes per vertex, so lazy Dijkstra is $1.1$ to $2.3$ times slower rather than
severely degraded (Hypothesis~1.3). And the pointer heaps make $1.2n$ to
$3.3n$ allocations in lazy mode, not the roughly $m$ predicted by
Hypothesis~3.2. A graph that really favours a cheap decrease-key needs edge
weights arranged so that relaxations keep succeeding, which uniform random
weights do not produce.

\subsection{Brodal-Okasaki's decrease-key is lazy Dijkstra with extra work}

Brodal-Okasaki was the only heap that ran faster in lazy mode, on every graph
and by up to $41\%$ ($9.57$ against $16.11$~ms on \code{uniform_dense}), and
it used half as much memory doing so ($1{,}920$ against $3{,}840$~KiB). The
implementation explains why. A functional heap has no handles, so
decrease-key cannot find and change an existing entry. Instead it inserts a
tombstone for the old (distance, vertex) pair into a negative queue and
inserts the new pair, and both stay in the buffer until they meet at the top
and cancel. Every decrease-key therefore costs two insertions now and two
delete-min operations later, where lazy Dijkstra pays one insertion and one
delete-min for the same relaxation. In effect, the \code{decrease_key} variant
is lazy Dijkstra with tombstones added, so it can only be slower. This also
overturns Hypothesis~3.1, which placed Brodal-Okasaki second-smallest in
memory: it has the largest peak in 13 of the 14 configurations.

\subsection{The heap with the fewest cache misses had the highest miss rate}

Our L1 miss-rate figure ranked the binary heap worst, which seemed to
contradict its speed. Counted per vertex, it has the fewest misses ($12.1$,
against $30.3$ for Brodal-Okasaki); its rate is high only because it makes so
few loads ($147$ per vertex against $763$). A rate divides by a denominator
that differs between implementations, so it can rank them in the opposite
order to their cost, which is why Section~4.3 reports counts instead.

\subsection{Our measurements were wrong twice before they were right}

The first hardware-counter results showed every heap running at an IPC close
to $1$ and retiring $26$ to $28$ million instructions on \code{smoke_1k}, even
though the heaps' runtimes on that graph ranged from $30$ to $230$
microseconds. The counters were recording the whole process, including
start-up and graph construction, rather than Dijkstra alone. With the counting
window fixed, the binary heap retires $0.40$ million instructions on that graph
and Brodal-Okasaki $1.24$ million. Later, a full rerun made the Fibonacci heap
about three times slower on several graphs while the other heaps were
unchanged. Dividing the \code{perf} cycle counts by the measured runtimes
showed that the processor had run at about $1.7$~GHz during those
benchmarks, against roughly $4.8$~GHz elsewhere: the laptop had throttled.
Both errors produced plausible-looking results, and both were caught only by
checking one measurement against another.

\subsection{What held}

Several predictions did hold. The Fibonacci heap never beat the pairing heap,
as Hypothesis~1.2 expected of its cascading cuts and four-pointer nodes, and it
suffers the most from lazy processing ($2.28\times$ slower on
\code{uniform_dense}). Brodal-Okasaki was the slowest heap on every graph,
supporting Hypothesis~1.4, although its slowdown is better described as a
constant-factor cost of this implementation than as a consequence of the data
structure itself. Once the counters were restricted to the Dijkstra loop, they
explained the ranking clearly: the pairing and Fibonacci heaps execute $1.4$
and $2.5$ times the binary heap's instructions at a similar IPC, while
Brodal-Okasaki executes $4.6$ times as many at an IPC of about $1.0$, held
back by frequent branch mispredictions. Overall, the experiments show that
for this workload, how often each operation is actually used matters as much
as its asymptotic cost, and memory layout and constant factors decide the
rest.
